\section{Introduction}
\label{sec:intro}

Most physical theories begin from assumed structure:
a spacetime manifold, a Lagrangian density, or a symmetry group.
In all such cases, the framework presupposes what it seeks to explain.

VED takes a different approach.
Starting from the single statement \textit{there is difference},
the framework derives time, space, motion, and persistence as necessary
consequences of structural asymmetry --- without importing any external
scaffolding.

The central object is the \textbf{causal log matrix} $C_{ij}$,
defined on a discrete network of nodes.
$C_{ij}$ captures not the total history of causal influence between nodes,
but the portion of that history that is selectively retained as persistent structure.
From $C_{ij}$, a single self-referential equation generates the entire chain:

\begin{keybox}
\begin{equation}
  \text{difference}
  \;\to\; C_{ij}
  \;\to\; \rho
  \;\to\; \tau
  \;\to\; d_{ij}
  \;\to\; J
  \;\to\; \omega
  \;\to\; L
  \label{eq:chain}
\end{equation}
\end{keybox}

\noindent
See Fig.~1 for a visual summary.

This framework is not presented as a complete predictive theory, but as a minimal structural derivation of physical organization.
It is a structural proposal: a demonstration that the minimum ingredients
required for physical structure are considerably smaller than those
assumed in current foundational approaches.
The aim is coherence, not closure.

\bigskip

\noindent\textbf{Organization.}
\Cref{sec:axiom} states the minimal axiom and traces its immediate consequences.
\Cref{sec:time} introduces the causal log and derives time.
\Cref{sec:space} derives spatial distance from the symmetric component of $C_{ij}$.
\Cref{sec:flow} develops flow and vorticity from the antisymmetric component.
\Cref{sec:closure} introduces closure as the minimal persistent structure.
\Cref{sec:phase} establishes the phase structure of the network.
\Cref{sec:discussion,sec:conclusion} discuss implications and open problems.
